% ==============================================================================
% CEBE Interdisciplinary Fellowship -- CV template
%
% Max 2 pages per person. Copy this file once per applicant, e.g.:
%   cv-pi.tex, cv-cosupervisor1.tex, cv-cosupervisor2.tex, cv-candidate.tex
% and compile each separately, then merge with main.tex into the single
% submission PDF (see README.md).
%
% Required minimum content per the call:
%   - Name, title and affiliation
%   - PhD year
%   - Short statement on previous research achievements and relevant
%     impacts (e.g. policy, industrial, or public arenas)
%   - Number (not a name list) of current and former PhD students,
%     postdocs, and master's students
%   - 10 publications: 5 most important (last 15 years) + 5 recent
%     project-relevant ones (last 5 years)
%   - Link to Scopus profile or ORCID
% ==============================================================================

% ==============================================================================
% Instantiated from templates/cv-template.tex on 2026-09-02 at Giuseppe's
% request. Name/title/affiliation below are pre-filled from AGENTS.md's
% "Team" section where known; everything else (PhD year, achievements,
% supervision counts, publications, ORCID/Scopus) is still a placeholder.
% ==============================================================================

\documentclass[12pt,a4paper]{article}

\usepackage[a4paper,margin=2.5cm]{geometry}
\usepackage[T1]{fontenc}
\usepackage[utf8]{inputenc}
\usepackage{mathptmx}
\usepackage{setspace}
\singlespacing
\usepackage{enumitem}
\usepackage{xcolor}
\usepackage[colorlinks=true,linkcolor=blue,urlcolor=blue]{hyperref}

\newcommand{\pagemark}[1]{\typeout{PAGEMARK: #1 -- page \thepage}}

% Expected-length annotation (matches proposal/main-template.tex). Set
% \templatenotesfalse to hide before the final submission build.
\newif\iftemplatenotes
\templatenotestrue
\newcommand{\lengthnote}[1]{%
  \iftemplatenotes
    {\small\itshape\color{gray!70!black} [Target length: #1]}\par\vspace{4pt}
  \fi
}

\pagenumbering{gobble} % CVs typically don't need page numbers in a merged submission; remove if you want them

\begin{document}
\pagemark{CV Giuseppe Abbiati (PI) -- start}
%\lengthnote{max 2 pages per person (per the call)}

\begin{center}
% Full name inferred from git/GitHub identity (giuseppe.abbiati@gmail.com,
% github.com/giuseppeabbiati1984) -- confirm/correct if needed.
{\Large\bfseries Giuseppe Abbiati}\par\vspace{2pt}
[Associate Professor in Structural Mechanics]\par
Dept. of Civil and Architectural Engineering (CAE), Aarhus University, Denmark
\end{center}

\vspace{0.3cm}
\noindent\textbf{ORCID:} 0000-0002-5048-8505, \textbf{Scopus ID:} 37064203300,\textbf{Scopus h-index:}  14, \textbf{Scopus citations:} 817, \textbf{Pure:} \url{https://pure.au.dk/portal/en/persons/abbiati@cae.au.dk}

\vspace{0.2cm}
\noindent\textbf{Education}
\begin{itemize}[nosep]
\item Ph.D. in Structural Engineering, University of Trento, Italy. Doctoral thesis title: “Dynamic substructuring of complex hybrid systems based on time integration, model-order reduction and dynamic identification techniques” (April 7, 2014).
\item M.Sc. in Civil Engineering, Polytechnic of Turin, Italy (December 10, 2009).
\item B.Sc. in Civil and Environmental Engineering, Polytechnic of Turin, Italy (October 21, 2008).
\end{itemize}

\vspace{0.2cm}
\noindent\textbf{Employments}
\begin{itemize}[nosep]
\item Associate Professor--Head of Section for Structural Engineering, Department of Engineering, Section of Civil and Architectural Engineering, University of Aarhus, Denmark (2023-now).
\item Tenure-Track Assistant Professor, Department of Engineering, Section of Civil and Architectural Engineering, University of Aarhus, Denmark (2019-2023).
\item Post-doctoral researcher, Chair of Structural Dynamics and Earthquake Engineering, Department of Civil, Environmental and Geomatic Engineering, ETH Zurich, Switzerland (2014-2018).
\end{itemize}

\noindent\textbf{Fellowships}\\
Recipient of Swiss Government Excellence Scholarship for a 1-year post-doctoral fellowship at the Department of Civil, Environmental and Geomatic Engineering of ETH Zurich, Switzerland, 2014, REF. 2014.0588.

\vspace{0.2cm}
\noindent\textbf{Research achievements and impact}

\noindent Giuseppe Abbiati's main research contributions relate to hybrid testing and digital twins for civil engineering structures (buildings, bridges, wind parks). His cumulative portfolio of funded research and innovation projects sums up to 2.9 MEUR, of which +2 MEUR are with industrial partners; +700 kEUR currently ongoing (all related to digital twins for the built environment).

%[Short statement on previous research achievements and relevant impacts, including e.g. policy, industrial, or public-arena impact. A few sentences, not a full narrative CV.]

\vspace{0.2cm}
\noindent\textbf{Supervision record}

\noindent Current: 5 PhD students, 2 postdocs, 7master's students.\\
Total: 7 PhD students (2 already graduated), 4 postdocs, 45 master's students. \\
Member of the assessment committee for 9 doctoral thesis.

\vspace{0.2cm}
\noindent\textbf{Teaching experience}
Course responsible for 15 ECTS in the MSc (finite-element analysis, experimental mechanics, machine learning for civil engineering) and 5 ECTS in the BSc (probability and statistics) of Civil and Architectural Engineering, Aarhus University.

\vspace{0.2cm}
\noindent\textbf{Selected research project}
\begin{enumerate}[nosep]
    \item Enhancing Hybrid Digital Twinning for the Preservation of Built Heritage (HERITWIN), 2026-2029, Marie Curie Staff Exchange Network, European Commission 123 kEUR (1.117 MEUR total budget)
    \item Laboratories as data factories for DT of floating wind turbines operational from Day 1 (LAB4TWIN)	2026-2029 CET Partnership Joint Call 2025, 334 kEUR (1.160 MEUR total budget)
    \item GIS-based platform for seismic resilience of urban communities (GURU), REF. 4339-00005B
	2025-2028 (36m)	EUREKA – Innovation Beyond Borders - 2024-23234/NP/Resilience Call 2024	298 kEUR (1.82 MEUR total budget)
    \item Cyber-Physical Sensing for Machine and Structures (CP-SENS), REF. 2081-00006B	2022-2027 (48m) Danish Innovation Fund 480 kEUR (2.55 MEUR total budget).
\end{enumerate}

\vspace{0.2cm}
\noindent\textbf{Selected publications}

\noindent\textit{Five most important (last 15 years):}
\begin{enumerate}[nosep]%[leftmargin=1.2cm]
  \item Scussolini, L., Abbiati, G., \& Ceravolo, R. (2026). Identification of Volterra time-varying parameters in dynamical systems under random excitation. Journal of Sound and Vibration, 645, 120128, DOI:10.1016/j.jsv.2026.120128.
  \item Pizarro, D., M. Kovarbašić, G. Abbiati, and B. Stojadinović, (2025). “Multi‐Axial Subassemblage Testing System for Hybrid Simulation with Large‐Scale Structural Components.” Earthquake Engineering \& Structural Dynamics, 54(8):2084-2105,DOI:10.1002/eqe.4351.
  \item Stamenov, D., Abbiati, G., \& Sauder, T. (2025). Numerical estimation of generalized frequency response functions from time series data using NARX. Mechanical Systems and Signal Processing, 239, 113278, DOI:10.1016/j.ymssp.2025.113278. 
  \item G. Abbiati, M. Broccardo, I. Abdallah, S. Marelli, F. Paolacci, (2021). Seismic Fragility Analysis based on Artificial Ground Motions and Surrogate Modeling of Validated Structural Simulators, Earthquake Engineering and Structural Dynamics, 50(9): 2314-2333, ISSN 00988847, DOI:10.1002/eqe.3448.
  \item G. Abbiati, V. La Salandra, O.S. Bursi, L. Caracoglia (2018). A Composite Experimental Dynamic Substructuring Method Based on Partitioned Algorithms and Localized Lagrange Multipliers. Mechanical Systems and Signal Processing, 100:85-112, ISSN 08883270, DOI:10.1016/j.ymssp.2017.07.020.
  \item Abbiati G., R. Ceravolo, C. Surace (2015). Time-dependent Estimators for On-Line Monitoring of Full-Scale Structures under Ambient Excitation, Mechanical Systems and Signal Processing, 60:166-181, ISSN 08883270, DOI:10.1016/j.ymssp.2014.10.018.
\end{enumerate}

\noindent\textit{Five recent, project-relevant (last 5 years):}
\begin{enumerate}[nosep]%[leftmargin=1.2cm]
  \item P. Talasila, D. Tcherniak, A. M.D. Jensen, S Mahato, J. V. Medom, M. D. Ulriksen, G. Abbiati, A. Schörghofer-Queiroz, P. G. Larsen, L. Damkilde, (2026). A Digital Twin Platform for Structural Health Monitoring, Computer-Aided Civil and Infrastructure Engineering, 50:100086, ISSN 1093-9687, DOI:10.1016/j.cacaie.2026.100086.
  \item Hansen, S. T., Thule, C., Gomes, C., Lausdahl, K. G., Madsen, F. P., Abbiati, G., \& Larsen, P. G. (2024). Co-simulation at different levels of expertise with Maestro2. Journal of Systems and Software, 209, 111905, DOI:10.1016/j.jss.2023.111905.
  \item Oakes, B, C Gomes, G Abbiati, E Kamburjan, E E. Baş, and S Engelsgaard. “Towards Ontological Service-Driven Engineering of Digital Twins.” In 1st International Conference on Engineering Digital Twins (EDTconf 2024). Linz, Austria: to appear, 2024, DOI:10.1145/3652620.3688261.
  \item Baş, E. E., Abbiati, G., Gonçalves Gomes, C. Â., Jassmann, U., \& Larsen, P. G. (2024). Digitalization of Large-Scale Testing Facilities for the Wind Industry: DIGIT-BENCH Digital Twin. Journal of Physics: Conference Series, 2767(4), 042033, DOI:10.1088/1742-6596/2767/4/042033.
  \item Abbiati, G., C. Gomes, M. Sandberg, Z. Kazemi, S. T. Hansen, and P. G. Larsen. “Modelling for Digital Twins.” In The Engineering of Digital Twins, edited by J. Fitzgerald, C. Gomes, and P. G. Larsen, 89–127. Cham: Springer International Publishing, 2024. DOI:10.1007/978-3-031-66719-0\textunderscore5.
\end{enumerate}

\pagemark{CV Giuseppe Abbiati (PI) -- end}
\end{document}
